\section{Information-Theoretic Analysis of Acknowledgments}
\label{app:info_theory}

To quantify why acknowledgment pruning under greedy decoding does not significantly degrade accuracy, we measure five information-theoretic properties of each message relative to its predecessor, grouped by communication function type. Metrics are computed on Qwen and Llama round-robin traces (300 tasks each); Gemma is excluded due to insufficient acknowledgment count ($n = 11$).

\begin{table}[ht]
\centering
\caption{Information-theoretic properties by function type (Qwen/GSM8K/RR, $N = 1{,}500$ messages). All ack vs.\ non-ack comparisons significant at $p < 10^{-10}$ (Mann--Whitney $U$) except where noted.}
\label{tab:info_theory_qwen}
\small
\begin{tabular}{@{}lrrrrr@{}}
\toprule
\textbf{Metric} & \textbf{Ack} & \textbf{Info Prov.} & \textbf{Reasoning} & \textbf{Verif.} & \textbf{Coord.} \\
\midrule
Bigram novelty & \textbf{0.057} & 0.317 & 0.749 & 0.434 & --- \\
Token length (median) & \textbf{4} & 95 & 117 & 53 & --- \\
Jaccard overlap & \textbf{0.686} & 0.653 & 0.313 & 0.482 & --- \\
Semantic similarity & \textbf{0.783} & 0.896 & 0.823 & 0.841 & --- \\
Information density & \textbf{0.971} & 0.565 & 0.566 & 0.685 & --- \\
\bottomrule
\end{tabular}
\end{table}

\begin{table}[ht]
\centering
\caption{Information-theoretic properties by function type (Llama/GSM8K/RR, $N = 1{,}470$ messages).}
\label{tab:info_theory_llama}
\small
\begin{tabular}{@{}lrrrrr@{}}
\toprule
\textbf{Metric} & \textbf{Ack} & \textbf{Info Prov.} & \textbf{Reasoning} & \textbf{Coord.} & \textbf{Verif.} \\
\midrule
Bigram novelty & 0.656 & 0.759 & 0.723 & 0.925 & 0.625 \\
Token length (mean) & 59.8 & 120.1 & 124.8 & 72.4 & 51.5 \\
Jaccard overlap & 0.342 & 0.332 & 0.336 & 0.191 & 0.283 \\
Semantic similarity & 0.636 & 0.817 & 0.763 & 0.453 & 0.706 \\
Information density & \textbf{0.774} & 0.472 & 0.519 & 0.741 & 0.718 \\
\bottomrule
\end{tabular}
\end{table}

\paragraph{Surface-level information content.}
Table~\ref{tab:surface_info} complements the above metrics with lexical diversity (type--token ratio, TTR) and bigram/trigram overlap with the immediately preceding message, computed over word-level tokens.
Across both backbones, acknowledgments exhibit the highest n-gram overlap with their predecessors (Qwen: 0.757 bigram, 0.733 trigram; Llama: 0.426, 0.363). This pattern confirms that acknowledgments primarily echo prior content.
Conversely, information provision and coordination messages show near-zero overlap ($< 0.11$ for Llama; $< 0.26$ for Qwen); these low overlaps indicate genuinely novel content.
This quantitative gap directly explains the perturbation-sensitivity dissociation: perturbing an acknowledgment alters surface form but not functional information, since the replaced content was redundant with the preceding turn.

\begin{table}[ht]
\centering
\caption{Surface-level information metrics by function type and backbone (GSM8K/round-robin). \textit{Bi-OL} and \textit{Tri-OL}: fraction of current-message bigrams/trigrams also present in the preceding message. TTR: type--token ratio (unique words / total words). Gemma excluded ($n_\text{ack} = 11$).}
\label{tab:surface_info}
\small
\begin{tabular}{@{}llrrrrr@{}}
\toprule
\textbf{Backbone} & \textbf{Function Type} & \textbf{$N$} & \textbf{Tokens} & \textbf{TTR} & \textbf{Bi-OL} & \textbf{Tri-OL} \\
\midrule
\multirow{4}{*}{Qwen}
  & Acknowledgment        & 951 & 39.2  & 0.820 & \textbf{0.757} & \textbf{0.733} \\
  & Verification          & 106 & 86.7  & 0.529 & 0.859 & 0.813 \\
  & Information provision  &  82 & 99.9  & 0.520 & 0.255 & 0.185 \\
  & Reasoning guidance    & 361 & 59.2  & 0.654 & 0.349 & 0.297 \\
\midrule
\multirow{5}{*}{Llama}
  & Acknowledgment        & 670 & 82.7  & 0.595 & \textbf{0.426} & \textbf{0.363} \\
  & Verification          & 112 & 115.8 & 0.422 & 0.420 & 0.369 \\
  & Information provision  & 192 & 62.1  & 0.665 & 0.235 & 0.184 \\
  & Reasoning guidance    & 376 & 69.2  & 0.624 & 0.168 & 0.105 \\
  & Coordination          & 120 & 37.4  & 0.845 & 0.105 & 0.064 \\
\bottomrule
\end{tabular}
\end{table}

\paragraph{Dual mechanism.} Qwen acknowledgments are informationally vacuous: median bigram novelty is 0\% and median Jaccard overlap with the preceding message is 1.0, meaning they are near-verbatim echoes. Pruning removes no unique information. Llama acknowledgments are longer and carry more novel content (bigram novelty 65.6\%), but their pruning benefit arises from a different mechanism: they inject errors at 14.3\% with 0\% correction, so removal cuts an error-propagation channel rather than removing redundancy.

\paragraph{Per-round acknowledgment trends.} The proportion of acknowledgments varies across conversation rounds and backbones. Qwen shows monotonic increase from 0\% (round~1) to 96\% (round~5), consistent with later rounds serving primarily confirmatory roles. Llama exhibits an inverted-U pattern peaking at 81\% in round~3 before declining; this decline suggests a transition from acknowledgment to substantive re-engagement in later rounds. Gemma produces near-zero acknowledgments across all rounds, consistent with its position-organized profile where the acknowledgment category is functionally absent.

\subsection{Heuristic Detector Performance}
\label{app:heuristic}

Backbone-specific heuristic detectors using only deterministic features (length, keywords, n-gram repetition) achieve F1\textsubscript{GSM8K} $= 0.845$ (Qwen, rule: character length $< 30$) and $0.805$ (Llama, composite rule: keyword matching $+$ length $< 300$ $+$ 4-gram repetition $> 0.3$) with 0--1\,pp pruning gap versus oracle labels, validated via 5-fold GroupKFold CV (grouped by task). Composite rules generalize cross-domain: Llama F1 $> 0.71$ on StratQA and ARC. Qwen's single-threshold rule degrades cross-domain (F1 $= 0.21$ on StratQA), motivating the multi-feature approach for deployment. Full development details appear in \S\ref{sec:heuristic}.
